Vision-language models (VLMs) that answer free-text questions about images
have progressed rapidly on general-domain benchmarks, driven by large-scale
image-text pretraining~\cite{radford2021clip} and efficient adaptation of
frozen encoders and language models~\cite{li2023blip2}. Applying these
models to medical imaging is attractive: a system that can answer a
clinician's question about a chest radiograph or CT slice could support
triage, education, and second-opinion workflows. However, general-domain
VLMs are trained overwhelmingly on natural images and web-scale captions,
and even biomedical-specific variants~\cite{li2023llavamed} are trained on
a fixed snapshot of literature and cannot cite the specific evidence behind
a given answer. An answer generated purely from a model's parameters is
difficult to audit: a clinician evaluating the system's output has no way
to inspect what, if anything, the model's answer was based on.

Retrieval-augmented generation (RAG) offers one route to addressing this.
Originally proposed for knowledge-intensive text tasks~\cite{lewis2020rag},
RAG couples a parametric generator with a non-parametric memory: instead of
answering from weights alone, the model is given retrieved supporting
evidence and conditions its answer on that evidence. In principle this
offers two benefits for medical VQA specifically: (1) the retrieved
evidence is inspectable, giving a clinician something concrete to check the
answer against, and (2) the model may produce more accurate answers by
grounding its response in similar, verified cases rather than relying
solely on what it memorized during pretraining.

Whether retrieval augmentation actually improves accuracy on medical VQA,
however, is an empirical question, not a foregone conclusion. Retrieved
evidence could just as easily distract a model, dilute its attention across
irrelevant content, or be ignored outright if the model was never trained
to make use of retrieved context. Answering this question rigorously
requires a controlled comparison: the same base model, the same test
questions, the same generation parameters, differing only in whether
retrieved evidence is present in the prompt.

\subsection{Contributions}

This paper makes three contributions:

\begin{enumerate}
  \item An open, reproducible pipeline for retrieval-augmented medical VQA
        built entirely on openly licensed, non-credentialed datasets ---
        ROCOv2~\cite{ruckert2024rocov2} as the retrieval corpus and
        VQA-RAD~\cite{lau2018vqarad} as the evaluation benchmark --- so
        the full pipeline can be reproduced without a data use agreement.
  \item A controlled empirical comparison of a vision-language model with
        and without retrieved evidence, using the same checkpoint,
        generation parameters, and scoring code for both conditions
        (Section~\ref{sec:method}), evaluated with both standard automatic
        metrics and paired statistical significance testing
        (Section~\ref{sec:experimental_setup}).
  \item A documented, config-driven, tested codebase intended to be
        extensible to larger or credentialed corpora (e.g., MIMIC-CXR) in
        future work, released alongside this paper.
\end{enumerate}
